\section{More Noetherian rings}
\label{section-Noetherian-again}


\begin{lemma}
\label{lemma-Noetherian-basic}
Let $R$ be a Noetherian ring.
Any finite $R$-module is of finite presentation.
Any submodule of a finite $R$-module is finite.
The ascending chain condition holds for $R$-submodules
of a finite $R$-module.
\end{lemma}

\begin{proof}
We first show that any submodule $N$ of a finite $R$-module
$M$ is finite. We do this by induction on the number of
generators of $M$. If this number is $1$, then $N = J/I \subset
M = R/I$ for some ideals $I \subset J \subset R$. Thus the definition
of Noetherian implies the result. If the number of generators of
$M$ is greater than $1$, then we can find a short exact sequence
$0 \to M' \to M \to M'' \to 0$ where $M'$ and $M''$ have fewer
generators. Note that setting $N' = M' \cap N$ and $N'' = \Im(N \to
M'')$ gives a similar short exact sequence for $N$. Hence the result
follows from the induction hypothesis
since the number of generators of $N$ is at most the number of
generators of $N'$ plus the number of generators of $N''$.

\medskip\noindent
To show that $M$ is finitely presented just apply the previous result
to the kernel of a presentation $R^n \to M$.

\medskip\noindent
It is well known and easy to prove that the ascending chain condition for
$R$-submodules of $M$ is equivalent to the condition that every submodule
of $M$ is a finite $R$-module. We omit the proof.
\end{proof}

\begin{lemma}[Artin-Rees]
\label{lemma-Artin-Rees}
Suppose that $R$ is Noetherian, $I \subset R$ an ideal.
Let $N \subset M$ be finite $R$-modules.
There exists a constant $c > 0$ such that
$I^n M \cap N  =  I^{n-c}(I^cM \cap N)$ for all $n \geq c$.
\end{lemma}

\begin{proof}
Consider the ring $S = R \oplus I \oplus I^2 \oplus \ldots
= \bigoplus_{n \geq 0} I^n$. Convention: $I^0 = R$.
Multiplication maps $I^n \times I^m$
into $I^{n + m}$ by multiplication in $R$.
Note that if $I = (f_1, \ldots, f_t)$
then $S$ is a quotient of the Noetherian ring $R[X_1, \ldots, X_t]$.
The map just sends the monomial $X_1^{e_1}\ldots X_t^{e_t}$
to $f_1^{e_1}\ldots f_t^{e_t}$. Thus $S$ is Noetherian.
Similarly, consider the module $M \oplus IM \oplus I^2M \oplus \ldots
= \bigoplus_{n \geq 0} I^nM$. This is a finitely generated $S$-module.
Namely, if $x_1, \ldots, x_r$ generate $M$ over $R$, then they also generate
$\bigoplus_{n \geq 0} I^nM$ over $S$. Next, consider the
submodule $\bigoplus_{n \geq 0} I^nM \cap N$.
This is an $S$-submodule, as is easily verified. By
Lemma \ref{lemma-Noetherian-basic} it is finitely generated as
an $S$-module,
say by $\xi_j \in \bigoplus_{n \geq 0} I^nM \cap N$, $j = 1, \ldots, s$.
We may assume by decomposing each $\xi_j$ into its homogeneous
pieces that each $\xi_j \in I^{d_j}M \cap N$ for some $d_j$.
Set $c = \max\{d_j\}$. Then for all $n \geq c$ every element
in $I^nM \cap N$ is of the form $\sum h_j \xi_j$ with
$h_j \in I^{n - d_j}$. The lemma now follows from this and the trivial
observation that $I^{n-d_j}(I^{d_j}M \cap N) \subset I^{n-c}(I^cM \cap N)$.
\end{proof}

\begin{lemma}
\label{lemma-map-AR}
Suppose that $0 \to K \to M \xrightarrow{f} N$ is an
exact sequence of finitely generated modules
over a Noetherian ring $R$. Let $I \subset R$ be an ideal.
Then there exists a $c$ such that
$$
f^{-1}(I^nN) = K + I^{n-c}f^{-1}(I^cN)
\quad\text{and}\quad
f(M) \cap I^nN \subset f(I^{n - c}M)
$$
for all $n \geq c$.
\end{lemma}

\begin{proof}
Apply Lemma \ref{lemma-Artin-Rees} to
$\Im(f) \subset N$ and note that
$f : I^{n-c}M \to I^{n-c}f(M)$ is surjective.
\end{proof}

\begin{lemma}[Krull's intersection theorem]
\label{lemma-intersect-powers-ideal-module-zero}
Let $R$ be a Noetherian local ring. Let $I \subset R$ be
a proper ideal. Let $M$ be a finite $R$-module.
Then $\bigcap_{n \geq 0} I^nM = 0$.
\end{lemma}

\begin{proof}
Let $N = \bigcap_{n \geq 0} I^nM$.
Then $N = I^nM \cap N$ for all $n \geq 0$.
By the Artin-Rees Lemma \ref{lemma-Artin-Rees}
we see that $N = I^nM \cap N \subset IN$ for
some suitably large $n$. By Nakayama's Lemma \ref{lemma-NAK}
we see that $N = 0$.
\end{proof}

\begin{lemma}
\label{lemma-intersection-powers-ideal-module}
Let $R$ be a Noetherian ring. Let $I \subset R$ be an ideal.
Let $M$ be a finite $R$-module. Let $N = \bigcap_n I^n M$.
\begin{enumerate}
\item For every prime $\mathfrak p$, $I \subset \mathfrak p$ there
exists a $f \in R$, $f \not \in \mathfrak p$ such that $N_f = 0$.
\item If $I$ is contained in the Jacobson radical
of $R$, then $N = 0$.
\end{enumerate}
\end{lemma}

\begin{proof}
Proof of (1). Let $x_1, \ldots, x_n$ be generators for the module $N$,
see Lemma \ref{lemma-Noetherian-basic}. For every prime
$\mathfrak p$, $I \subset \mathfrak p$ we see that
the image of $N$ in the localization $M_{\mathfrak p}$
is zero, by Lemma \ref{lemma-intersect-powers-ideal-module-zero}.
Hence we can find $g_i \in R$, $g_i \not \in \mathfrak p$
such that $x_i$ maps to zero in $N_{g_i}$. Thus
$N_{g_1g_2\ldots g_n} = 0$.

\medskip\noindent
Part (2) follows from (1) and Lemma \ref{lemma-characterize-zero-local}.
\end{proof}

\begin{remark}
\label{remark-intersection-powers-ideal}
Lemma \ref{lemma-intersect-powers-ideal-module-zero} in particular implies
that $\bigcap_n I^n = (0)$ when $I \subset R$ is a non-unit ideal in a
Noetherian local ring $R$. More generally, let $R$ be a Noetherian ring and
$I \subset R$ an ideal. Suppose that $f \in \bigcap_{n \in \mathbf{N}} I^n$.
Then Lemma \ref{lemma-intersection-powers-ideal-module}
says that for every prime ideal $I \subset \mathfrak p$
there exists a $g \in R$, $g \not \in \mathfrak p$
such that $f$ maps to zero in $R_g$. In algebraic geometry we
express this by saying that ``$f$ is zero in an open neighbourhood
of the closed set $V(I)$ of $\Spec(R)$''.
\end{remark}

\begin{lemma}[Artin-Tate]
\label{lemma-Artin-Tate}
Let $R$ be a Noetherian ring. Let $S$ be a finitely
generated $R$-algebra. If $T \subset S$ is an $R$-subalgebra such
that $S$ is finitely generated as a $T$-module, then $T$ is of
finite type over $R$.
\end{lemma}

\begin{proof}
Choose elements $x_1, \ldots, x_n \in S$ which generate $S$ as an $R$-algebra.
Choose $y_1, \ldots, y_m$ in $S$ which generate $S$ as a $T$-module.
Thus there exist $a_{ij} \in T$ such that
$x_i = \sum a_{ij} y_j$. There also exist $b_{ijk} \in T$ such
that $y_i y_j = \sum b_{ijk} y_k$. Let $T' \subset T$ be the
sub $R$-algebra generated by $a_{ij}$ and $b_{ijk}$. This is a finitely
generated $R$-algebra, hence Noetherian. Consider the algebra
$$
S' = T'[Y_1, \ldots, Y_m]/(Y_i Y_j - \sum b_{ijk} Y_k).
$$
Note that $S'$ is finite over $T'$, namely as a $T'$-module it is
generated by the classes of $1, Y_1, \ldots, Y_m$.
Consider the $T'$-algebra homomorphism $S' \to S$ which maps
$Y_i$ to $y_i$. Because $a_{ij} \in T'$ we see that $x_j$ is
in the image of this map. Thus $S' \to S$ is surjective.
Therefore $S$ is finite over $T'$ as well. Since $T'$ is Noetherian
we conclude that $T \subset S$ is finite over $T'$ and
we win.
\end{proof}






















